\documentclass[12pt]{article}

\pagestyle{empty}
\setcounter{secnumdepth}{2}

\topmargin=0cm
\oddsidemargin=0cm
\textheight=22.0cm
\textwidth=16cm
\parindent=0cm
\parskip=0.15cm
\topskip=0truecm
\raggedbottom
\abovedisplayskip=3mm
\belowdisplayskip=3mm
\abovedisplayshortskip=0mm
\belowdisplayshortskip=2mm
\normalbaselineskip=12pt
\normalbaselines

\begin{document}

\vspace*{0.5in}
\centerline{\bf\Large Test Document}

\vspace*{0.5in}
\centerline{\bf\Large Team PJA}

\vspace*{0.5in}
\centerline{\bf\Large April 7, 2019}

\vspace*{1.5in}
\begin{table}[htbp]
\caption{Team}
\begin{center}
\begin{tabular}{|r | c|}
\hline
Name & ID Number \\
\hline\hline
David Boivin & 40004941\\
\hline
X & Y\\
\hline
\end{tabular}
\end{center}
\end{table}

\clearpage

\section{Introduction}

This document describes the testing procedures planned and executed on the Codenames project by team {\it PJA} and their results. The tests described have been carefully planned in the early phases of developement (iteration 1) and have been gradually written and executed as the project neare completion and the necessary portions of code were written.
\\

\section{Test Plan}

The Test Plan is split into 3 sections, System Level tests, Subsystem Level tests, and Unit Tests. Each of these sections represent a level of integration of the program and have been scheduled in order of integration level, namely Unit tests, Subsystem tests and finally System tests.
\\

{\bf System Level: } Tests the integration and cohesion of all subsystems to form the finished project. This section also includes usability tests which tests if our program can be interacted with as the requirements specify and make sure said interactions produce the correct output efficiently. These tests are of utmost importance as they regulate whether the program functions according to specification and as such were the first to be conceived yet the last to be executed since all subsystems had to be complete and operational prior. Test execution took place in iteration 3 of development and consist of running the program and testing different situations and user interactions, such as card selection and appropriate reactions based on strategy, control button reactions.
\\

{\bf Subsystem Level: } Tests the cooperation between low level units to form a subsystem, which in turn performs a major and complex operation for the system as a whole. Such operation inlcude use of different strategies, creation of the program's internal model and controller interaction with model. These tests are designed to ensure that subsystems perform their designated task correctly in preperation for system wide testing. This level of testing was done by using the Facade testing pattern providing a sandbox to test components together. Such tests were planned after the system tests as per the white to black box testing scheme and executed in iteration 2 of developement due to the requirement of functioning units.
\\

{\bf Unit Test: } Tests the lowest level functions of each Unit of code, or class within the code, which make up the components of a subsystem. This level of testing is necessarily designed and executed very early in developement (iteration 1) to alleviate concerns of flaws and/or bugs prior to subsystem testing and full blown system testing. It is important to note that not all classes have been thoroughly tested, only important data classes and linking classes within the model and control portions of the program.
\\


\subsection{System Level Test Cases}

\subsubsection{Test Case 1: Player Card Selection} \label{tc:1}

\noindent
{\bf Purpose}\\
Tests the ability for a user to select a card on their turn

\noindent
{\bf Input Specification}\\
Create a new game with a User and an AI (Random or Sequential). If it is an AI start, then press Next Move Control Button until it is User's turn. Click on a card.

\noindent
{\bf Expected Output}\\
\emph{Case 1:} User Team Card\\
It is the User Team's turn to play again.\\
\emph{Case 2:} AI Team Card or Bystander\\
It is the AI's turn to play.\\
\emph{Case 3:} Assassine Card\\
The game ends.\\
\emph{Case 4:} Selected Card \\
Nothing happens.

\noindent
{\bf Traces to Use Cases}\\
The User must be able to play the game as an Operative.


\subsubsection{Test Case 2: Next Move Control Button} \label{tc:2}

\noindent
{\bf Purpose}\\
Tests the ability of the AI to make a move.

\noindent
{\bf Input Specification}\\
Create a AI vs. AI new game (Random/Sequential). If Red start, press Next Move until Blue team turn. Press Next Move control button.

\noindent
{\bf Expected Output}\\
\emph{Case 1:} Blue AI 1 Team Card\\
It is the AI 1's turn to play again.\\
\emph{Case 2:} AI Team Card or Bystander\\
It is the AI 2's turn to play.\\
\emph{Case 3:} Assassine Card\\
The game ends.\\

\noindent
{\bf Traces to Use Cases}\\
The User can play against AI with intellegent Random or Iterative strategies. 


\subsubsection{Test Case 3: Undo Control Button} \label{tc:3}

\noindent
{\bf Purpose}\\
Tests the ability of the user to undo any number of moves.

\noindent
{\bf Input Specification}\\
Create a new game. In any order, press Next Move control button twice and select 2 cards. Press Undo control button.

\noindent
{\bf Expected Output}\\
The game state is exactly what it was before playing the 4th move.

\noindent
{\bf Traces to Use Cases}\\
Part of the Undo/Redo system

\subsubsection{Test Case 4: Redo Control Button } \label{tc:4}

\noindent
{\bf Purpose}\\
Tests the ability of the user to redo moves, if there are any.

\noindent
{\bf Input Specification}\\
Create a new game. In any order, press Next Move control button twice and select 2 cards. Press Undo control button twice. Press Redo control button.

\noindent
{\bf Expected Output}\\
The game state is exactly what it was before playing the 4th move.

\noindent
{\bf Traces to Use Cases}\\
Part of the Undo/Redo system.

\subsubsection{Test Case 5: Keycard Control Button} \label{tc:5}

\noindent
{\bf Purpose}\\
Tests the ability of the user to see the keycard.

\noindent
{\bf Input Specification}\\
Start a new game. Press the Keycard control button. Find the Black assassin card on Keycard screen. Close Keycard screen. Press the card where the assassin should be according to the keycard.

\noindent
{\bf Expected Output}\\
Assasin card should be flipped and game should end.

\noindent
{\bf Traces to Use Cases}\\
User should be able to see the Keycard.

\subsubsection{Test Case 6: Quit Control Button} \label{tc:6}

\noindent
{\bf Purpose}\\
Tests whether the game can be closed.

\noindent
{\bf Input Specification}\\
Start a new game. Press Close control button.

\noindent
{\bf Expected Output}\\
Program should end.

\noindent
{\bf Traces to Use Cases}\\
Easily close the program.

\subsubsection{Test Case 7: New Game Control Button} \label{tc:7}

\noindent
{\bf Purpose}\\
Tests the ability of the user to start a new game.

\noindent
{\bf Input Specification}\\
Start application, press New Game control button. Select AI types (Sequential/Random).Press Continue button.

\noindent
{\bf Expected Output}\\
Game should be different.

\noindent
{\bf Traces to Use Cases}\\
Create a new game.

\subsection{Subsystem Level Test Cases}

{\it
All test cases for testing at the subsystem level.
}

{\it
One subsection per subsystem
}

\subsubsection{Subsystem 1: Hint Subsystem}

\subsubsection{Subsystem 2: Model Creation Subsystem}

\subsubsection{Subsystem 3: Strategy Subsystem}

\subsubsection{Subsystem 4: Point Subsystem}

\subsection{Unit Test cases}

{\it
All test cases for testing at the unit level.
}

{\it
One subsection per unit
}

\subsubsection{Unit X}

\section{Test Results}

{\it
List the tests, indicating which passed and which did not pass.
List requirements indicating the percentage of tests that passed for that requirement.
}

\section{References}

\subsection{Description of Input Files}

\subsubsection{Appendix 1: CodenamesDatabase.txt}

This input file is filled with the words used inside the program and has a format:
\begin{center}
\{word: {\it word1} hints: {\it hint11, hint12, hint13, ...} \}\\
\{word: {\it word2} hints: {\it hint21, hint22, hint23, ...} \}\\...
\end{center}
The file contains 100 words where each word has at least 3 hints and each hint relates to a minimum of 2 words. Hints are not considered words, though a word can be used as a hint.

The CodenamesDatabase.txt can be found at: \\
\begin{center}
COMP354PJA/client/src/CodenameDatabase.txt
\end{center}

\section{Description of Output Files}

Describe/include test expected output that are output files.
\\\\
------------\- Abysnth: We Probably dont need this section

\end{document}
